\documentclass[11pt,a4paper]{article}

% Compile with: pdflatex sthyra-lead-capture-guide.tex
% Run pdflatex twice if you want the table of contents page numbers updated.

\usepackage[T1]{fontenc}
\usepackage[utf8]{inputenc}
\usepackage{lmodern}
\usepackage[margin=22mm]{geometry}
\usepackage{xcolor}
\usepackage{hyperref}
\usepackage{bookmark}
\usepackage{enumitem}
\usepackage{tabularx}
\usepackage{longtable}
\usepackage{booktabs}
\usepackage{array}
\usepackage{amssymb}
\usepackage{listings}
\usepackage{fancyhdr}
\usepackage{titlesec}
\usepackage{microtype}

\definecolor{SthyraGreen}{HTML}{287B63}
\definecolor{SthyraMint}{HTML}{42B995}
\definecolor{Ink}{HTML}{18201D}
\definecolor{Muted}{HTML}{5F6965}
\definecolor{Rule}{HTML}{DCE4E0}
\definecolor{CodeBg}{HTML}{F5F8F7}
\definecolor{WarningBg}{HTML}{FFF8E8}

\hypersetup{
  colorlinks=true,
  linkcolor=SthyraGreen,
  urlcolor=SthyraGreen,
  citecolor=SthyraGreen,
  pdftitle={Sthyra CRM Lead Capture and Integration Guide},
  pdfauthor={Sthyra CRM}
}

\pagestyle{fancy}
\fancyhf{}
\lhead{\textcolor{Muted}{STHYRA CRM}}
\rhead{\textcolor{Muted}{Lead Capture Guide}}
\cfoot{\thepage}
\renewcommand{\headrulewidth}{0.4pt}
\renewcommand{\headrule}{\hbox to\headwidth{\color{Rule}\leaders\hrule height \headrulewidth\hfill}}

\titleformat{\section}{\Large\bfseries\color{Ink}}{\thesection}{0.7em}{}
\titleformat{\subsection}{\large\bfseries\color{Ink}}{\thesubsection}{0.7em}{}
\titleformat{\subsubsection}{\normalsize\bfseries\color{Ink}}{\thesubsubsection}{0.7em}{}
\setlist[itemize]{leftmargin=*,itemsep=3pt,topsep=5pt}
\setlist[enumerate]{leftmargin=*,itemsep=5pt,topsep=5pt}
\setlength{\parindent}{0pt}
\setlength{\parskip}{6pt}

\lstdefinestyle{sthyra}{
  backgroundcolor=\color{CodeBg},
  basicstyle=\ttfamily\small,
  breaklines=true,
  breakatwhitespace=false,
  columns=fullflexible,
  frame=single,
  rulecolor=\color{Rule},
  frameround=tttt,
  showstringspaces=false,
  tabsize=2,
  xleftmargin=4pt,
  xrightmargin=4pt,
  aboveskip=8pt,
  belowskip=8pt
}
\lstset{style=sthyra}

\newcommand{\product}{Sthyra CRM}
\newcommand{\crmurl}{https://crm.sthyra.com}
\newcommand{\menu}[1]{\textbf{#1}}
\newcommand{\code}[1]{\texttt{#1}}
\newcommand{\placeholder}[1]{\texttt{<#1>}}
\newcommand{\warningbox}[1]{%
  \vspace{4pt}%
  \noindent\colorbox{WarningBg}{%
    \parbox{\dimexpr\linewidth-2\fboxsep\relax}{\textbf{Important:} #1}%
  }%
  \vspace{4pt}%
}

\begin{document}

\begin{titlepage}
  \color{Ink}
  \vspace*{22mm}
  {\Large\bfseries\textcolor{SthyraGreen}{STHYRA CRM}\par}
  \vspace{18mm}
  {\Huge\bfseries Lead Capture and Integration Guide\par}
  \vspace{5mm}
  {\Large Adding leads from files, websites, advertising platforms, partners and APIs\par}
  \vspace{16mm}
  \textcolor{Rule}{\rule{\textwidth}{1pt}}
  \vspace{8mm}

  \begin{tabular}{@{}ll@{}}
    Document version: & 1.0 \\
    Product version: & Sthyra CRM V1 \\
    Production URL: & \url{https://crm.sthyra.com} \\
    Last updated: & 30 September 2026 \\
  \end{tabular}

  \vfill
  {\small\textcolor{Muted}{For workspace administrators, marketing teams, website developers and integration partners.}\par}
\end{titlepage}

\tableofcontents
\newpage

\section{Purpose and Scope}

This guide explains how leads enter \product, how to configure each capture method inside the CRM, and what must be configured on the source website or external platform. It covers the lead-capture features currently available in V1:

\begin{itemize}
  \item validated CSV import;
  \item first-party website and landing-page forms;
  \item Google Ads lead-form webhooks;
  \item generic webhooks for partners, portals and automation tools;
  \item the trusted single and bulk lead-intake APIs; and
  \item optional Google Ads conversion feedback after a lead has entered the CRM.
\end{itemize}

\warningbox{The current V1 Leads page does not expose a manual ``Create lead'' form. For a one-off lead, use a one-row CSV import or a configured capture endpoint. Do not tell users to look for an Add Lead button that is not present in this build.}

\subsection{Choose the correct method}

\begin{tabularx}{\textwidth}{@{}>{\bfseries}p{31mm}X X@{}}
\toprule
Method & Best for & Authentication \\
\midrule
CSV import & Initial migration, historical data and controlled batch uploads & Signed-in user with Data Import permission \\
Website form & A first-party website or campaign landing page & Allowed browser origin; no secret exposed in the browser \\
Google lead form & Google Ads lead-form assets & Google webhook key generated by Sthyra \\
Partner/portal webhook & Channel partners, property portals, Zapier, Make or another server & One-time server submission secret \\
Lead-intake API & Trusted backend systems and high-control integrations & Global bearer secret managed by the Sthyra operator \\
Google Ads conversions & Sending qualified/site-visit/booking/won outcomes back to Google & Separate Google Data Manager OAuth connection \\
\bottomrule
\end{tabularx}

\section{Before Adding Leads}

Complete these steps once for each workspace or project.

\subsection{Confirm access}

The person configuring acquisition should be a company administrator with access to every relevant project. CSV operators need the \code{DATA\_IMPORT} permission. Users who only need to work incoming records need the appropriate lead-view and lead-update permissions.

\subsection{Create or verify the project}

In \menu{Settings $>$ Administration $>$ Projects}, confirm that the destination project is active. Record its numeric \code{project\_id}; CSV and API integrations require this value.

Every active project must have an active initial lead stage. Review this under \menu{Settings $>$ Current project $>$ Lead stages}.

\subsection{Configure attribution}

\begin{enumerate}
  \item Open \menu{Settings $>$ Administration $>$ Lead sources}.
  \item Create a source such as Website, Google Ads, Channel Partner or Property Portal.
  \item Open \menu{Settings $>$ Administration $>$ Campaigns}.
  \item Create the campaign and, where appropriate, associate it with its lead source.
  \item Keep the generated source and campaign identifiers for CSV/API use.
\end{enumerate}

Campaign and source identifiers are UUIDs. A campaign selected for a lead must belong to the selected source when the campaign has a source association.

\subsection{Configure assignment and response rules}

Select the project in Settings and open \menu{Lead configuration}.

\begin{itemize}
  \item Select a default source and campaign if required.
  \item Choose Manual, Round robin or Load balanced assignment.
  \item Enable \menu{Automatically assign incoming leads} if new leads should be routed immediately.
  \item Set the response SLA.
  \item Configure queues and routing rules under the current project when automatic assignment is enabled.
\end{itemize}

\section{Method 1: Import Leads from CSV}

CSV import is the safest option for migrations and controlled batches. The CRM validates the entire file before writing anything and imports a maximum of 1,000 rows per operation. If one committed row fails, the whole batch is rolled back.

\subsection{Steps inside Sthyra CRM}

\begin{enumerate}
  \item Open \menu{Settings $>$ Administration $>$ Data tools}.
  \item Under \menu{Import CSV}, choose \menu{Leads} as the record type.
  \item Upload a CSV file or paste CSV content into the editor.
  \item Click \menu{Validate}. Resolve every reported row and column error.
  \item Click \menu{Import data} only after validation succeeds.
  \item Open \menu{Leads}, select the destination project and confirm the imported records.
\end{enumerate}

\subsection{Required and optional columns}

Each lead row requires:

\begin{itemize}
  \item \code{project\_id};
  \item \code{first\_name}; and
  \item at least one of \code{email} or \code{phone\_number}.
\end{itemize}

Supported optional columns include:

\begin{lstlisting}
last_name,email,phone_number,alternate_phone_number,date_of_birth,is_nri,
country,company_works_at,is_married,anniversary_date,address,source_id,
campaign_id,sub_source,temperature,customer_type,preferred_location,
preferred_config,preferred_facing,preferred_floor,preferred_view,budget,
buying_reason
\end{lstlisting}

The \code{temperature} value, when supplied, must be \code{cold}, \code{warm} or \code{hot}. The budget must be a non-negative number. UUID values must reference active records in the same workspace.

\subsection{Example CSV}

\begin{lstlisting}
project_id,first_name,last_name,email,phone_number,temperature,budget,sub_source
42,Aarav,Shah,aarav@example.com,+919876543210,hot,8500000,September launch
42,Meera,Rao,meera@example.com,,warm,6200000,Referral drive
\end{lstlisting}

During import, Sthyra reuses an existing contact when its email or phone matches; otherwise it creates a new contact. Each valid row then creates a lead in the project's initial stage and applies the project's routing rules.

\section{Method 2: Capture Leads from Your Website}

Use this method for a website, microsite or landing page controlled by your organization.

\subsection{Steps inside Sthyra CRM}

\begin{enumerate}
  \item Open \menu{Marketing $>$ Forms \& webhooks}.
  \item Click \menu{New endpoint}.
  \item Enter a clear endpoint name.
  \item Choose \menu{Website form}.
  \item Select the destination project, lead source and optional campaign.
  \item Add every allowed website origin, one per line, for example \url{https://www.example.com}. Use only exact origins; do not include a path.
  \item Create the endpoint.
  \item Copy the public form key and generated script. Save the one-time server secret in a password manager even if the browser integration will not use it.
\end{enumerate}

Allowed origins may also use a subdomain wildcard such as \code{*.example.com}. Prefer exact production and staging origins wherever possible.

\subsection{Steps on the website}

Add the generated script and mark the HTML form with \code{data-sthyra-form}. Replace \placeholder{PUBLIC-FORM-KEY} with the key shown by Sthyra.

\begin{lstlisting}
<form id="enquiry-form" data-sthyra-form>
  <label>
    Full name
    <input name="full_name" required>
  </label>
  <label>
    Email
    <input name="email" type="email">
  </label>
  <label>
    Phone
    <input name="phone_number" type="tel">
  </label>
  <label>
    Preferred location
    <input name="preferred_location">
  </label>
  <label>
    Configuration
    <input name="preferred_config">
  </label>
  <label>
    Reason for buying
    <textarea name="buying_reason"></textarea>
  </label>
  <button type="submit">Request a callback</button>
</form>

<script
  src="https://crm.sthyra.com/sthyra-marketing.js"
  data-form-key="<PUBLIC-FORM-KEY>"
  defer>
</script>
\end{lstlisting}

The script automatically captures the landing page, referrer, UTM parameters and supported advertising click identifiers such as \code{gclid}, \code{gbraid}, \code{wbraid} and \code{fbclid}.

Supported canonical website field names are:

\begin{lstlisting}
full_name, first_name, last_name, email, phone_number, country,
company_works_at, address, preferred_location, preferred_config,
budget, buying_reason
\end{lstlisting}

At least \code{first\_name} or \code{full\_name}, plus either \code{email} or \code{phone\_number}, must be submitted.

\subsection{Custom JavaScript submission}

For a React application or a custom form handler, call the script API directly:

\begin{lstlisting}
const result = await window.SthyraMarketing.submit(
  {
    first_name: "Aarav",
    last_name: "Shah",
    email: "aarav@example.com",
    phone_number: "+919876543210",
    preferred_location: "Kolkata",
    preferred_config: "3 BHK",
    budget: 8500000
  },
  {
    adUserDataConsent: "granted",
    adPersonalizationConsent: "denied"
  }
);
console.log(result.lead_id);
\end{lstlisting}

Use a JavaScript number for \code{budget}. Values collected automatically through standard HTML \code{FormData} are strings, so convert the budget before using the custom API if it must be stored as a number.

The HTML form emits \code{sthyra:submitting}, \code{sthyra:submitted} and \code{sthyra:error} events. Use them to display loading, success and error messages without replacing the CRM integration.

\section{Method 3: Google Ads Lead Forms}

This method receives leads created through a Google Ads lead-form asset. It is separate from the Google Data Manager conversion integration described later.

\subsection{Steps inside Sthyra CRM}

\begin{enumerate}
  \item Open \menu{Marketing $>$ Forms \& webhooks} and click \menu{New endpoint}.
  \item Select \menu{Google lead form} as the provider.
  \item Select the project, source and campaign.
  \item Enter the Google lead-form asset ID exactly as shown in Google Ads.
  \item Create the endpoint.
  \item Copy both the webhook URL and the one-time \menu{Google key}.
\end{enumerate}

The production webhook URL has this form:

\begin{lstlisting}
https://crm.sthyra.com/api/marketing/webhooks/google/<PUBLIC-FORM-KEY>
\end{lstlisting}

\subsection{Steps in Google Ads}

\begin{enumerate}
  \item Sign in to the Google Ads account that owns the campaign.
  \item Open the lead-form asset used by the campaign.
  \item Find its webhook or lead-delivery configuration.
  \item Paste the Sthyra webhook URL into the webhook URL field.
  \item Paste the Sthyra Google key into Google's key field.
  \item Save the asset and send a test lead.
  \item In Sthyra, confirm the event under \menu{Marketing $>$ Touchpoints}, then confirm the resulting record under \menu{Leads}.
\end{enumerate}

Sthyra recognizes Google's full name, first name, last name, email, phone number and city fields, together with optional property-specific fields named \code{PREFERRED\_LOCATION}, \code{PROPERTY\_TYPE} and \code{PURCHASE\_TIMELINE}.

If Google reports an ``unknown form'' error, compare the form ID sent by Google with the provider form ID saved in Sthyra. They must match exactly.

\section{Method 4: Partner, Portal or Generic Webhook}

Use this method for a channel partner, property portal, automation platform or another backend service.

\subsection{Steps inside Sthyra CRM}

\begin{enumerate}
  \item Open \menu{Marketing $>$ Forms \& webhooks}.
  \item Click \menu{New endpoint}.
  \item Select \menu{Generic webhook}, \menu{Channel partner} or \menu{Property portal}.
  \item Select the destination project, source and campaign.
  \item Create the endpoint.
  \item Copy the endpoint URL and the one-time server submission secret.
  \item Store the secret in the external platform's secret manager. It cannot be displayed again.
\end{enumerate}

\subsection{Steps in the external platform}

Configure an HTTP \code{POST} request to the endpoint supplied by Sthyra. Set:

\begin{lstlisting}
Content-Type: application/json
X-Sthyra-Form-Secret: <SERVER-SUBMISSION-SECRET>
\end{lstlisting}

Send a stable, unique \code{event\_id}. Retrying the same event ID is safe and does not create a second intake event.

\begin{lstlisting}
curl -X POST \
  'https://crm.sthyra.com/api/marketing/forms/<PUBLIC-FORM-KEY>/submit' \
  -H 'Content-Type: application/json' \
  -H 'X-Sthyra-Form-Secret: <SERVER-SUBMISSION-SECRET>' \
  -d '{
    "event_id": "partner-acme-2026-000184",
    "fields": {
      "first_name": "Meera",
      "last_name": "Rao",
      "email": "meera@example.com",
      "phone_number": "+919812345678",
      "preferred_location": "New Town",
      "preferred_config": "2 BHK"
    },
    "attribution": {
      "event_type": "partner_lead",
      "channel": "partner",
      "medium": "referral",
      "external_campaign_id": "partner-september-2026"
    }
  }'
\end{lstlisting}

If an external platform cannot send a custom HTTP header, do not place the secret in a URL or browser script. Send the platform webhook first to a secure middleware function, then forward it to Sthyra with the required header.

\section{Method 5: Trusted Lead-Intake API}

The lead-intake API is intended for Sthyra-managed backend integrations. It uses the server-side \code{LEAD\_INTAKE\_SECRET}, which is global and must not be distributed to browsers, customers or untrusted partners.

\subsection{Single lead}

\begin{lstlisting}
curl -X POST 'https://crm.sthyra.com/api/lead-intake' \
  -H 'Authorization: Bearer <LEAD-INTAKE-SECRET>' \
  -H 'Content-Type: application/json' \
  -d '{
    "idempotency_key": "erp:lead:93842",
    "contact": {
      "first_name": "Aarav",
      "last_name": "Shah",
      "email": "aarav@example.com",
      "phone_number": "+919876543210"
    },
    "lead": {
      "project_id": 42,
      "source_id": "<SOURCE-UUID>",
      "campaign_id": "<CAMPAIGN-UUID>",
      "sub_source": "ERP sync",
      "temperature": "warm",
      "budget": 8500000
    }
  }'
\end{lstlisting}

Instead of a new \code{contact} object, the request may contain an existing \code{contact\_id}. Supply one or the other, never both.

\subsection{Bulk lead intake}

Send between 1 and 100 envelopes per request:

\begin{lstlisting}
curl -X POST 'https://crm.sthyra.com/api/lead-intake/bulk' \
  -H 'Authorization: Bearer <LEAD-INTAKE-SECRET>' \
  -H 'Content-Type: application/json' \
  -d '{
    "events": [
      {
        "idempotency_key": "partner:batch-77:row-1",
        "contact": {
          "first_name": "Meera",
          "email": "meera@example.com"
        },
        "lead": {"project_id": 42, "temperature": "hot"}
      }
    ]
  }'
\end{lstlisting}

The single endpoint returns HTTP 201 when processing completes and HTTP 202 when the event is accepted but needs further handling. Reusing the same company and idempotency key returns the existing event with \code{duplicate: true}. The bulk endpoint returns HTTP 207 with a result for every input item.

\section{Google Ads Conversion Feedback}

Lead capture brings a person into the CRM. Conversion feedback sends later CRM outcomes back to Google Ads. Configure this only when the acquisition team wants Google to optimize against qualified leads, completed site visits, confirmed bookings or won opportunities.

\subsection{Google-side preparation}

\begin{enumerate}
  \item Create the required conversion actions in Google Ads/Data Manager.
  \item Record the operating Google Ads account ID and, if applicable, the manager/login account ID.
  \item Record the numeric conversion action ID for each milestone used by the workspace.
  \item Configure a Google OAuth web client with the Data Manager scope and the callback below.
\end{enumerate}

\begin{lstlisting}
https://crm.sthyra.com/api/marketing/google/callback
\end{lstlisting}

The deployment must contain \code{GOOGLE\_MARKETING\_CLIENT\_ID}, \code{GOOGLE\_MARKETING\_CLIENT\_SECRET}, \code{GOOGLE\_MARKETING\_REDIRECT\_URI} and \code{INTEGRATION\_TOKEN\_ENCRYPTION\_KEY}.

\subsection{Sthyra-side setup}

\begin{enumerate}
  \item Open \menu{Marketing $>$ Integrations}.
  \item Click \menu{Add Google Ads}.
  \item Enter an integration name, operating account ID and optional manager/login account ID.
  \item Enter the conversion action IDs for the milestones that should be exported.
  \item Save the integration and click \menu{Connect Google}.
  \item Authorize the Google account that can access the selected Ads account.
  \item Review delivery status under \menu{Marketing $>$ Conversions}. Retry failed jobs after correcting the reported issue.
\end{enumerate}

\section{Verification After Any Setup}

Always submit one clearly identifiable test lead before declaring an integration live.

\begin{enumerate}
  \item Confirm that the request reports success.
  \item Open \menu{Marketing $>$ Touchpoints} and find the test event.
  \item Open \menu{Leads}, select the correct project and find the contact.
  \item Verify source, campaign, sub-source, initial stage and attribution.
  \item Confirm that routing assigned the correct owner or queue.
  \item Confirm that SLA tracking began when applicable.
  \item Delete or mark the test record invalid according to workspace policy.
\end{enumerate}

\section{Troubleshooting}

\begin{longtable}{@{}p{34mm}p{48mm}p{\dimexpr\textwidth-90mm\relax}@{}}
\toprule
Symptom & Likely cause & Resolution \\
\midrule
\endhead
\code{401 Unauthorized} & Missing or incorrect webhook/API secret & Use the current secret in the correct header. If it was rotated, update the external platform. \\
\addlinespace
\code{403 Origin is not allowed} & Browser origin is absent from the endpoint allowlist & Add the exact scheme and hostname, such as \url{https://www.example.com}, then retry from that site. \\
\addlinespace
\code{404 Form not found} & Wrong public key, paused endpoint or incorrect URL & Copy the endpoint again and confirm that it is marked Live in Sthyra. \\
\addlinespace
\code{422 Invalid submission} & Missing first name/contact method, invalid field type or inactive reference & Inspect the response, correct the payload and resubmit with the same event ID only when retrying the same lead. \\
\addlinespace
\code{422 Unknown Google form} & Google form ID does not equal the ID stored in Sthyra & Update the capture endpoint or the Google lead-form configuration. \\
\addlinespace
\code{429 Too many requests} & Endpoint rate limit exceeded & Back off and retry later; investigate accidental loops or duplicate webhook deliveries. \\
\addlinespace
CSV validates with errors & Unknown header, invalid UUID, invalid phone/email or inaccessible project & Use only supported headers and IDs from the same workspace. Correct every reported row before committing. \\
\addlinespace
Touchpoint exists but lead is absent & Intake validation was quarantined or failed & Check the submitted contact and lead values, especially the project and required contact fields, then replay through the authorized operational process. \\
\addlinespace
Lead is unassigned & Auto-assignment is off or no routing rule matched & Review Lead configuration, queues, user availability and routing rules for the project. \\
\bottomrule
\end{longtable}

\section{Security and Data-Quality Rules}

\begin{itemize}
  \item Never place the server submission secret, lead-intake bearer secret or OAuth client secret in browser code, a public repository or a URL.
  \item Use HTTPS for every production source and callback.
  \item Store one-time credentials immediately; rotate them when exposed or when an integration owner changes.
  \item Use a unique idempotency key for each real lead event and reuse it only for retries of that same event.
  \item Request only the personal data needed for the sales process and show the appropriate consent notice on the source form.
  \item Keep source, campaign and project mappings explicit. Do not use a generic endpoint for unrelated campaigns when accurate attribution matters.
  \item Monitor failed submissions, unassigned leads and stale follow-ups during launch testing.
\end{itemize}

\section{Go-Live Checklist}

\begin{itemize}[label=$\square$]
  \item Destination project is active and has an initial lead stage.
  \item Lead source and campaign are active and correctly associated.
  \item Assignment strategy, queues, routing rules and SLA are configured.
  \item Endpoint provider and project mapping are correct.
  \item Production origin, webhook URL or OAuth callback is exact.
  \item Secrets are stored only in approved secret managers.
  \item Test lead appears in both Marketing Touchpoints and Leads.
  \item Source, campaign, owner, stage and contact details are correct.
  \item Duplicate retries do not create another intake event.
  \item Error monitoring and an integration owner are assigned.
\end{itemize}

\section{Reference URLs}

\begin{itemize}
  \item Sthyra CRM: \url{https://crm.sthyra.com}
  \item Website script: \url{https://crm.sthyra.com/sthyra-marketing.js}
  \item Google Ads lead-form webhook guidance: \url{https://developers.google.com/google-ads/webhook/docs/overview}
  \item Google Data Manager API: \url{https://developers.google.com/data-manager/api}
\end{itemize}

\vfill
\begin{center}
  \textcolor{Muted}{End of guide}
\end{center}

\end{document}
